% !PRINTED 319
\weberpage{319}
%<<CONT
are of the 8th degree, which shall not be entered into more closely
here (cf. Bd. I, \S.\,185).

\begin{center}
\S.\ 84.

\gesperrt{Biquadratic Equations.}
\end{center}

We have restricted ourselves to the consideration of the primitive
metacyclic equations of the 8th degree, because the imprimitive ones
come back to equations of the 4th degree. In order, then, to be able to
represent the roots of the latter also in the manner of the preceding
section, we have still to seek an analogous representation of the roots
of biquadratic equations. Here the same considerations as in
\S.\,82 are applicable, only in an essentially simpler form.

We denote the roots of the biquadratic equation by
\weq{1}{X_{0,0}, X_{1,0}, X_{0,1}, X_{1,1},}

and can then represent the whole permutation group of the 24th
degree as a binary linear congruence group for the modulus 2:
\weq{2}{P = SQ.}

The homogeneous group $S$ contained therein is of the 6th
degree and consists of the substitutions
\weq{3}{\sigma = \begin{pmatrix}1, & 0\\0, & 1\end{pmatrix},
\begin{pmatrix}1, & 0\\1, & 1\end{pmatrix},
\begin{pmatrix}0, & 1\\1, & 0\end{pmatrix},
\begin{pmatrix}0, & 1\\1, & 1\end{pmatrix},
\begin{pmatrix}1, & 1\\1, & 0\end{pmatrix},
\begin{pmatrix}1, & 1\\0, & 1\end{pmatrix},}

and one may regard
\weq{4}{\sigma_1 = \begin{pmatrix}0, & 1\\1, & 0\end{pmatrix}, \quad
\sigma_2 = \begin{pmatrix}1, & 1\\0, & 1\end{pmatrix}}

as the generating substitutions of the group $S$.

We now define, as in \S.\,80, the three quantities:
\weq{5}{\begin{aligned}
\Psi_{1,0} &= X_{0,0} - X_{1,0} + X_{0,1} - X_{1,1}\\
\Psi_{0,1} &= X_{0,0} + X_{1,0} - X_{0,1} - X_{1,1}\\
\Psi_{1,1} &= X_{0,0} - X_{1,0} - X_{0.1} + X_{1,1}, % [sic] printed with a period after 0 instead of a comma
\end{aligned}}

so that $\Psi^2_{1,0}, \Psi^2_{0,1}, \Psi^2_{1,1}$, but also the product
$\Psi_{1,0}\Psi_{0,1}\Psi_{1,1}$, remain unchanged by the substitutions
of the cyclic group $Q$.
